% =============================================================================
% Kapitel 6: Fazit und Handlungsempfehlungen
% =============================================================================

\section{Fazit und Handlungsempfehlungen}
\label{sec:fazit}

Die Fallstudie hat fünf Failure Cases des ROSI-Systems auf Prozess- und
Betriebssystemebene untersucht und damit die drei eingangs formulierten
Forschungsfragen beantwortet.

Zu FF2 (Systembeobachtbarkeit) zeigt Case~1 die größte Lücke: Drei von vier
Service-Prozessen fallen ohne Erkennung aus. Da das System über kein externes Alerting
verfügt, bleibt die Time-to-Detection in diesen Fällen unbegrenzt. Das erhöht die
Wiederherstellungszeit und senkt nach Gleichung~\ref{eq:availability} die Verfügbarkeit.
Zu FF1 (Ressourcenverwaltung) belegt Case~2 die fehlende Rückfallebene des flüchtigen
Bildspeichers, während Case~5 eine angenommene Race Condition widerlegt. Zu FF3
(Datenpersistenz und Kaskaden) zeigen Case~3 und Case~4 zwei unterschiedliche Mechanismen:
einen systemweiten Stillstand durch Backpressure sowie ein durch die
Entwicklungskonfiguration maskiertes Verbindungsrisiko.

\subsection{Übergreifende Befunde}
\label{sec:fazit-uebergreifend}

Zwei Muster ziehen sich durch die Ergebnisse. Erstens fehlt eine durchgängige
\textit{Graceful Degradation}: Ein lokaler Fehler -- ein abgestürzter Prozess, ein
fehlgeschlagener Schreibvorgang, eine tote Verbindung -- führt zum Ausfall des
Gesamtsystems statt zu einem kontrollierten Teilbetrieb. Zweitens bestehen
\textit{blinde Flecken im Monitoring}: Weder Queue-Tiefen noch Speicherfüllstände noch der
Zustand der Service-Prozesse werden zur Laufzeit überwacht. Besonders Case~3 zeigt, dass
eine verlangsamte Hilfskomponente über eine gemeinsame Queue die eigentliche Pipeline
blockieren kann.

Beide Muster wirken über die MTTR auf die Verfügbarkeit. Ohne automatische Erkennung
(Abschnitt~\ref{sec:problemstellung}) bleibt die Time-to-Detection an die Anwesenheit von
Personal gebunden und kann im 24/7-Betrieb ohne Bereitschaftsdienst Stunden betragen.
Nach Gleichung~\ref{eq:availability} senkt jede Verlängerung der MTTR die Verfügbarkeit
unmittelbar -- unabhängig davon, wie selten der Fehler auftritt. Eine automatische
Erkennung, die die Time-to-Detection von Stunden auf Sekunden verkürzt, verbessert die
Verfügbarkeit damit selbst dann deutlich, wenn die Ausfallrate unverändert bleibt. Genau
hier setzen die folgenden Empfehlungen an.

\subsection{Priorisierte Handlungsempfehlungen}
\label{sec:fazit-empfehlungen}

Tabelle~\ref{tab:empfehlungen} fasst die abgeleiteten Maßnahmen zusammen. Die ersten vier
Maßnahmen verbessern direkt die Verfügbarkeit: Sie verkürzen entweder die
Time-to-Detection oder entkoppeln einen lokalen Fehler vom Gesamtsystem. Case~5 erfordert
keinen Sofortfix, eignet sich aber als Anlass für Step-5-Monitoring.

\begin{table}[htbp]
    \centering
    \caption{Priorisierte Handlungsempfehlungen}
    \label{tab:empfehlungen}
    \begin{tabularx}{\textwidth}{llX}
        \hline
        \textbf{Case} & \textbf{Aufwand} & \textbf{Maßnahme} \\
        \hline
        1 & gering   & Supervision aller Service-Prozesse oder \texttt{systemd}-Unit \\
        2 & gering   & Fehlerbehandlung in der Bildspeicherung + Füllstands-Monitoring \\
        3 & minimal  & \texttt{put(timeout)} mit Drop-on-Full in der Log-Queue \\
        4 & minimal  & \texttt{close\_old\_connections()} bzw. \texttt{CONN\_MAX\_AGE = 0} \\
        5 & --       & kein Sofortfix; optional Referenzzähler und Step-5-Monitoring \\
        \hline
    \end{tabularx}
\end{table}

\subsection{Ausblick}
\label{sec:fazit-ausblick}

Die wichtigsten Hebel für die Produktionsreife liegen in einer durchgängigen
Prozessüberwachung sowie einem Monitoring der bislang unbeobachteten Ressourcen. Beide
Maßnahmen adressieren nicht nur einzelne Cases, sondern die übergreifenden Muster
fehlender Beobachtbarkeit und fehlender Graceful Degradation. Eine
Übertragung der Untersuchung auf das spätere Kundensystem -- mit abweichender Hardware und
aktiver Produktionskonfiguration -- bleibt offen und würde insbesondere das in Case~4
analytisch hergeleitete Verbindungsrisiko empirisch absichern.
